\documentclass[11pt,a4paper]{article}
\usepackage[top=1.0in, bottom=1.0in, left=1.1in, right=1.1in]{geometry}
\usepackage{graphicx}
\usepackage[sort&compress, super, numbers]{natbib}
\usepackage[export]{adjustbox}
\usepackage[ngerman]{babel}
\usepackage[utf8]{inputenc}
\usepackage[T1]{fontenc}

% novelty statement, a cover letter, and a title page including an authorship statement, a data statement, the number of words in abstract/main text and the number of figure/table/boxes.
% recommended and opposed reviewers, and recommended and opposed editorial board members with whom they are in conflict. 
\begin{document}

\pagenumbering{gobble}

\noindent \includegraphics[width=0.4\textwidth, right]{letterhead/ubc-logo-2018-fullsig-blue-cmyk.png}

\noindent Dear Dr. Cottingham,\\

\noindent Please consider our paper, `Current environments and evolutionary history shape forest temporal assembly' for publication as an article in \emph{Ecology}. %We present results from large scale experiments of spring phenology in forest communities from across North America. Our work illustrates the importance of species evolutionary history in shaping spring phenology and the consistency of phenological cues across diverse species and geographic gradients.
\vspace{1.5 ex}\\
\noindent The timing of plant life history events—such as budburst and leafout—have shifted with climate change, with increasing concern for how these changes will impact communities and how they assemble in time. Species timings are, however, are highly variable, especially in temperate systems where the effects of climate change are predicted to be greatest.  We performed a large-scale experiment of leaftout with 47 forest species and over 1500 plants in order to decompose this variability into predictable responses to environmental cues,  allowing us to make inferences at the population, species and community-levels. In using a phylogenetic Bayesian model, we were also able to test the importance of evolutionary relationships in shaping species cure responses. We found that systematic differences in species' timings, which were able to decompose into stable responses to temperature and daylength and attributed to the degree of species' shared evolutionary histories. But species responses to leafout cues were surprisingly consistent across populations. Collectively our findings further our understanding the factors that shape forest communities and their temporal assembly. 
\vspace{1.5ex}\\
\noindent
Two reviewers highlighted the strengths of our analysis, having found the approach we took `elegant' and that our work contributes `important insights' to the field. Their comments, questions and suggested clarifications have greatly improved the manuscript. In response to the reviewer comments, we have provided greater detail in our methods and more explicitly defined several key terms throughout the manuscript to improve clarity. We also performed an additional analysis that further supports our key findings and their interpretation. We have also revised the discussion to provide a more in-depth overview of the literature and a more detailed explanation of the implications and importance of our findings to forecasting and predicting changes in forest communities with climate change. Finally, we revised our figures to improve their clarity and interpretability. 
\vspace{1.5ex}\\
\noindent We believe the new manuscript has been greatly improved, as discussed in our point-by-point responses. The manuscript is 23 pages or 4866 words with a 201 word abstract, and 5 figures. Our manuscript is not under consideration elsewhere. We hope you find it suitable for publication in \emph{Ecology}, and look forward to hearing from you. 
%\noindent We recommend the following reviewers: Dr. Meredith Zettlemoyer, Dr. Mason Heberling, Dr. Jason Fridley, Dr. Rong Yu, and Dr. Ameila Caffara. 
\vspace{1ex}\\
\noindent Sincerely, \\
\includegraphics[scale=.4]{letterhead/shot.png} \\ 
\noindent Deirdre Loughnan\\
\noindent Sentinels of Change Postdoctoral Fellow\\ 
\noindent Hakai Institute $|$ UBC Department of Zoology\\
\newpage
\vspace{-5ex}
\bibliographystyle{refs/bibstyles/nature.bst}% 

\bibliography{refs/traitors_mar23.bib}
\newpage


\end{document}